\unitchapter{1}{9}{People, Development \& Environment}{p1-09}

\unitmeta{Expected questions: \textasciitilde5 (10 marks) · Target: 4--5/5}

\begin{sylbox}

\begin{itemize}
\tightlist
\item[$\square$]
  Development and environment: Millennium Development and Sustainable Development Goals
\item[$\square$]
  Human and environment interaction: anthropogenic activities and their impacts on environment
\item[$\square$]
  Environmental issues: local, regional and global; air, water, soil, noise pollution; waste (solid, liquid, biomedical, hazardous, electronic); climate change and its socio-economic and political dimensions
\item[$\square$]
  Impacts of pollutants on human health
\item[$\square$]
  Natural and energy resources: solar, wind, soil, hydro, geothermal, biomass, nuclear and forests
\item[$\square$]
  Natural hazards and disasters: mitigation strategies
\item[$\square$]
  Environmental Protection Act (1986), NAPCC, international agreements/efforts: Montreal Protocol, Rio Summit, CBD, Kyoto Protocol, Paris Agreement, ISA
\end{itemize}

\end{sylbox}

\needspace{125pt}

\hypertarget{p1-09--1-mdgs--sdgs}{%
\section{MDGs → SDGs}\label{p1-09--1-mdgs--sdgs}}

Development was long measured by income alone. The global goals reflect a broader view: development must
also reduce poverty, improve health and education, and do so without exhausting the environment on which
future generations depend.

\begin{figure}[!htbp]
\centering
\begin{adjustbox}{max width=\linewidth}
\begin{tikzpicture}
  \draw[line width=5pt, fslate] (0,0) -- (12,0);
  \foreach \x/\y/\t in {0.4/2000/{Millennium Development Goals: 8 goals, 21 targets}, 6/2015/{Sustainable Development Goals: 17 goals, 169 targets — Agenda 2030}, 11.6/2030/{target year of the SDGs}} {
    \filldraw[fill=fwine, draw=white, line width=1pt] (\x,0) circle (5pt);
    \node[font=\small\bfseries, above=3mm] at (\x,0) {\y};
    \node[font=\footnotesize, text width=38mm, align=center, below=3mm] at (\x,0) {\t};
  }
\end{tikzpicture}
\end{adjustbox}
\caption{From the Millennium Development Goals to the Sustainable Development Goals.}
\end{figure}


\textbf{8 MDGs}: (1) eradicate extreme poverty \& hunger (2) universal primary education (3) gender equality
(4) reduce child mortality (5) improve maternal health (6) combat HIV/AIDS, malaria (7) environmental sustainability (8) global partnership.

\textbf{17 SDGs} (memorise numbers --- directly asked):

\begingroup\small

\begin{longtable}[]{@{}
  >{\raggedright\arraybackslash}p{(\columnwidth - 6\tabcolsep) * \real{0.0357}}
  >{\raggedright\arraybackslash}p{(\columnwidth - 6\tabcolsep) * \real{0.3631}}
  >{\raggedright\arraybackslash}p{(\columnwidth - 6\tabcolsep) * \real{0.0417}}
  >{\raggedright\arraybackslash}p{(\columnwidth - 6\tabcolsep) * \real{0.3661}}@{}}
\toprule\noalign{}
\begin{minipage}[b]{\linewidth}\raggedright
\textbf{\#}
\end{minipage} & \begin{minipage}[b]{\linewidth}\raggedright
\textbf{Goal}
\end{minipage} & \begin{minipage}[b]{\linewidth}\raggedright
\textbf{\#}
\end{minipage} & \begin{minipage}[b]{\linewidth}\raggedright
\textbf{Goal}
\end{minipage} \\
\midrule\noalign{}
\endhead
\bottomrule\noalign{}
\endlastfoot
1 & No Poverty & 10 & Reduced Inequalities \\
2 & Zero Hunger & 11 & Sustainable Cities \\
3 & Good Health \& Well-being & 12 & Responsible Consumption \& Production \\
\textbf{4} & \textbf{Quality Education} & \textbf{13} & \textbf{Climate Action} \\
5 & Gender Equality & 14 & Life Below Water \\
6 & Clean Water \& Sanitation & 15 & Life on Land \\
7 & Affordable \& Clean Energy & 16 & Peace, Justice \& Strong Institutions \\
8 & Decent Work \& Economic Growth & 17 & Partnerships for the Goals \\
9 & Industry, Innovation \& Infrastructure & & \\
\end{longtable}

\endgroup

5 Ps of SDGs: \textbf{People, Planet, Prosperity, Peace, Partnership}. In India, \textbf{NITI Aayog} monitors (SDG India Index, 2018).

\textbf{Sustainable development} (Brundtland Report "Our Common Future", 1987): \emph{development that meets the needs of the present without compromising the ability of future generations to meet their own needs.}

\needspace{125pt}

\hypertarget{p1-09--2-pollution--sources--health-effects}{%
\section{Pollution --- Sources \& Health Effects}\label{p1-09--2-pollution--sources--health-effects}}

Pollution is the addition of substances or energy to the environment faster than it can disperse or
neutralise them. Questions often pair a pollutant with the disease or disaster it caused, so learn the
pairs in the figure and table together.

\begin{figure}[!htbp]
\centering
\begin{adjustbox}{max width=\linewidth}
\begin{tikzpicture}
  \foreach \x/\col/\head/\items in {
      0/fslate/Air/{PM\textsubscript{2.5}, PM\textsubscript{10} — lungs, heart\\SO\textsubscript{2}, NO\textsubscript{x} — acid rain\\CO — binds haemoglobin\\ground-level O\textsubscript{3}\\lead — neurotoxic},
      1/fsage/Water/{fluoride — fluorosis\\arsenic — skin lesions\\mercury — Minamata\\cadmium — itai-itai\\nitrate — blue-baby},
      2/fsand/Soil/{pesticides, fertilisers\\heavy metals\\salinisation\\\phantom{x}\\\phantom{x}},
      3/frose/Noise/{above 85 dB harmful\\hearing loss\\stress, sleep loss\\\phantom{x}\\\phantom{x}}} {
    \node[box, fill=\col, minimum width=36mm] at (\x*3.95,0) {\textbf{\head}};
    \node[plain, text width=34mm, align=left, font=\footnotesize, anchor=north] at (\x*3.95,-0.6) {\items};
  }
\end{tikzpicture}
\end{adjustbox}
\caption{Major forms of pollution and their best-known health effects.}
\end{figure}


\begingroup\small

\begin{longtable}[]{@{}
  >{\raggedright\arraybackslash}p{(\columnwidth - 2\tabcolsep) * \real{0.3867}}
  >{\raggedright\arraybackslash}p{(\columnwidth - 2\tabcolsep) * \real{0.3750}}@{}}
\toprule\noalign{}
\begin{minipage}[b]{\linewidth}\raggedright
\textbf{Disease/\allowbreak{}Event}
\end{minipage} & \begin{minipage}[b]{\linewidth}\raggedright
\textbf{Cause}
\end{minipage} \\
\midrule\noalign{}
\endhead
\bottomrule\noalign{}
\endlastfoot
Minamata disease & Mercury (Japan) \\
Itai-itai & Cadmium (Japan) \\
Blue baby syndrome (methaemoglobinaemia) & Nitrates in water \\
Black-foot disease & Arsenic \\
Fluorosis & Excess fluoride \\
Bhopal gas tragedy (1984) & Methyl isocyanate (MIC) \\
Chernobyl (1986) & Nuclear accident (Ukraine) \\
Silicosis & Silica dust \\
Asbestosis & Asbestos \\
London smog (1952) & SO₂ + smoke (classical/\allowbreak{}sulphurous smog) \\
Photochemical smog & NOx + VOC + sunlight → O₃, PAN \\
\end{longtable}

\endgroup

\begin{itemize}
\tightlist
\item
  \textbf{NAAQS} (CPCB): PM2.5 annual 40 µg/m³, 24-hr 60; PM10 annual 60, 24-hr 100. \textbf{AQI} has 8 pollutants (PM10, PM2.5, NO₂, SO₂, CO, O₃, NH₃, Pb), 6 categories.
\item
  \textbf{WHO 2021} guideline PM2.5 annual: 5 µg/m³.
\item
  Greenhouse gases: CO₂, CH₄, N₂O, HFCs, PFCs, SF₆ (+ NF₃ under Doha amendment). GWP: CO₂=1, CH₄≈28, N₂O≈265, SF₆≈23,500.
\item
  Ozone depletion: \textbf{CFCs} (chlorine radicals); ozone measured in \textbf{Dobson Units}; ozone layer in stratosphere.
\end{itemize}

\needspace{160pt}

\hypertarget{p1-09--3-waste-management}{%
\section{Waste Management}\label{p1-09--3-waste-management}}

India\textquotesingle s 2016 set of waste rules shifted responsibility to the generators of waste and to producers
(extended producer responsibility). The guiding principle is the waste hierarchy: the best waste is the
waste never produced.

\begingroup\small

\begin{longtable}[]{@{}
  >{\raggedright\arraybackslash}p{(\columnwidth - 2\tabcolsep) * \real{0.1222}}
  >{\raggedright\arraybackslash}p{(\columnwidth - 2\tabcolsep) * \real{0.8778}}@{}}
\toprule\noalign{}
\begin{minipage}[b]{\linewidth}\raggedright
\textbf{Waste}
\end{minipage} & \begin{minipage}[b]{\linewidth}\raggedright
\textbf{Rules in India}
\end{minipage} \\
\midrule\noalign{}
\endhead
\bottomrule\noalign{}
\endlastfoot
Solid waste & Solid Waste Management Rules 2016 --- segregation at source (wet, dry, hazardous) \\
Plastic & Plastic Waste Management Rules 2016 (amended 2021/22) --- single-use plastic ban from 1 July 2022 \\
E-waste & E-Waste (Management) Rules 2016 → 2022 (EPR --- Extended Producer Responsibility) \\
Bio-medical & Bio-Medical Waste Management Rules 2016 --- colour-coded bins (yellow, red, white, blue) \\
Hazardous & Hazardous \& Other Wastes Rules 2016 \\
C\&D waste & Construction \& Demolition Waste Rules 2016 \\
\end{longtable}

\endgroup

Waste hierarchy (most → least preferred): \textbf{Prevent → Reduce → Reuse → Recycle → Recover → Dispose}.

\needspace{125pt}

\hypertarget{p1-09--4-natural--energy-resources}{%
\section{Natural \& Energy Resources}\label{p1-09--4-natural--energy-resources}}

India\textquotesingle s energy transition --- from coal towards solar, wind and other renewables --- is a frequent source of
questions, both for its targets and for its geography.

\begin{figure}[!htbp]
\centering
\begin{adjustbox}{max width=\linewidth}
\begin{tikzpicture}
  \node[grey, minimum width=34mm] (R) at (0,0) {\textbf{Energy resources}};
  \node[sage, text width=56mm] (RE) at (-3.4,-1.6) {\textbf{Renewable}\\\footnotesize solar, wind, hydro, biomass, geothermal, tidal};
  \node[rose, text width=56mm] (NR) at (3.4,-1.6) {\textbf{Non-renewable}\\\footnotesize coal, petroleum, natural gas, nuclear};
  \draw[arr] (R) -- (RE); \draw[arr] (R) -- (NR);
\end{tikzpicture}
\end{adjustbox}
\caption{Renewable resources are replenished by nature within a human lifetime; non-renewable ones are not.}
\end{figure}


\begin{itemize}
\tightlist
\item
  India\textquotesingle s targets: \textbf{500 GW non-fossil capacity by 2030}; \textbf{Net zero by 2070} (Panchamrit, COP26 Glasgow).
\item
  Largest share of India\textquotesingle s electricity: \textbf{coal (thermal)}. Largest renewable: \textbf{solar} (overtook wind).
\item
  Top states: solar --- Rajasthan, Gujarat; wind --- Gujarat, Tamil Nadu.
\item
  Nuclear: Tarapur (first, 1969). Geothermal sites: Puga Valley (Ladakh). Tidal potential: Gulf of Khambhat, Kutch.
\item
  India\textquotesingle s forest cover \textasciitilde21.7\% (ISFR 2021); tree + forest \textasciitilde25\%. National Forest Policy 1988 target: \textbf{33\%}.
\end{itemize}

\needspace{125pt}

\hypertarget{p1-09--5-natural-hazards--disasters}{%
\section{Natural Hazards \& Disasters}\label{p1-09--5-natural-hazards--disasters}}

A hazard becomes a disaster only when it strikes people who are vulnerable and unprepared. Disaster
management therefore works continuously, not only after the event.

\begin{figure}[!htbp]
\centering
\begin{tikzpicture}
  \def\R{2.3}
  \foreach \a/\n/\t/\s/\col in {90/M/Mitigation/{reduce risk before it strikes}/fsage, 0/P/Preparedness/{plans, drills, warnings}/fslate,
      -90/R/Response/{rescue, relief}/frose, 180/C/Recovery/{rehabilitate, rebuild better}/fsand} {
    \node[box, fill=\col, text width=27mm] (\n) at (\a:\R) {\textbf{\t}\\\footnotesize\s};
  }
  \draw[arr] (M) to[bend left=30] (P); \draw[arr] (P) to[bend left=30] (R);
  \draw[arr] (R) to[bend left=30] (C); \draw[arr] (C) to[bend left=30] (M);
  \node[capt, align=center] at (0,0) {disaster\\management\\cycle};
  \draw[decorate, decoration={zigzag, segment length=5pt, amplitude=2pt}, fwine, line width=0.8pt] (2.3,-2.6) -- (2.9,-1.9);
  \node[lbl, text=fwine, anchor=west] at (2.95,-1.85) {disaster strikes};
\end{tikzpicture}
\caption{The disaster management cycle. Mitigation and preparedness come \emph{before} the event; response
and recovery follow it.}
\end{figure}


\begin{itemize}
\tightlist
\item
  \textbf{Disaster Management Act 2005} → \textbf{NDMA} (chaired by PM), SDMA (CM), DDMA (Collector). \textbf{NDRF} (response force).
\item
  \textbf{Sendai Framework (2015--2030)} --- 4 priorities, 7 targets (succeeded Hyogo Framework 2005--2015).
\item
  Hazards: earthquake (Richter/moment magnitude), cyclone, flood, drought, landslide, tsunami (2004 Indian Ocean).
\item
  Risk = Hazard × Vulnerability / Capacity.
\end{itemize}

\needspace{196pt}

\hypertarget{p1-09--6-environmental-laws--international-agreements}{%
\section{Environmental Laws \& International Agreements}\label{p1-09--6-environmental-laws--international-agreements}}

Environmental law in India grew in response to specific crises --- the Bhopal gas tragedy led directly to
the Environment (Protection) Act, 1986. Internationally, cooperation advanced through a series of
conferences and protocols whose years and subjects are a staple of the examination.

\needspace{130pt}

\hypertarget{p1-09--indian-acts}{%
\subsection{Indian Acts}\label{p1-09--indian-acts}}

\begingroup\small

\begin{longtable}[]{@{}
  >{\raggedright\arraybackslash}p{(\columnwidth - 2\tabcolsep) * \real{0.4072}}
  >{\raggedright\arraybackslash}p{(\columnwidth - 2\tabcolsep) * \real{0.3044}}@{}}
\toprule\noalign{}
\begin{minipage}[b]{\linewidth}\raggedright
\textbf{Act}
\end{minipage} & \begin{minipage}[b]{\linewidth}\raggedright
\textbf{Year}
\end{minipage} \\
\midrule\noalign{}
\endhead
\bottomrule\noalign{}
\endlastfoot
Wildlife Protection Act & 1972 \\
Water (Prevention \& Control of Pollution) Act & 1974 (created CPCB) \\
Forest (Conservation) Act & 1980 \\
Air (Prevention \& Control of Pollution) Act & 1981 \\
\textbf{Environment (Protection) Act} & \textbf{1986} (post-Bhopal; umbrella act) \\
Biological Diversity Act & 2002 \\
National Green Tribunal Act & 2010 \\
\end{longtable}

\endgroup

\textbf{NAPCC (2008)} --- 8 National Missions: Solar, Enhanced Energy Efficiency, Sustainable Habitat, Water, Sustaining Himalayan Ecosystem, Green India, Sustainable Agriculture, Strategic Knowledge on Climate Change.

\needspace{95pt}

\hypertarget{p1-09--international}{%
\subsection{International}\label{p1-09--international}}

\begin{figure}[!htbp]
\centering
\begin{adjustbox}{max width=\linewidth}
\begin{tikzpicture}[y=-0.68cm]
  \draw[line width=1.6pt, fslate] (0,0.3) -- (0,10.7);
  \foreach \y/\yr/\ev [count=\k] in {1/1972/{Stockholm Conference — UNEP created}, 2/1985/{Vienna Convention for the ozone layer},
      3/1987/{Montreal Protocol on ozone-depleting substances; Brundtland Report}, 4/1992/{Rio Earth Summit — UNFCCC, CBD, Agenda 21},
      5/1997/{Kyoto Protocol — binding targets for developed countries (in force 2005)}, 6/2002/{Johannesburg Summit (Rio+10)},
      7/2012/{Rio+20 — "The Future We Want"}, 8/2015/{Paris Agreement (COP21): well below 2 °C, NDCs; International Solar Alliance},
      9/2016/{Kigali Amendment — phase-down of HFCs}, 10/2021/{COP26 Glasgow — India's net-zero target of 2070}} {
    \filldraw[fill=fwine, draw=white, line width=0.8pt] (0,\y) circle (3.5pt);
    \node[anchor=east, font=\small\bfseries] at (-0.3,\y) {\yr};
    \node[anchor=west, font=\small, text width=125mm] at (0.3,\y) {\ev};
  }
\end{tikzpicture}
\end{adjustbox}
\caption{Milestones in global environmental cooperation.}
\end{figure}


\begingroup\small

\begin{longtable}[]{@{}
  >{\raggedright\arraybackslash}p{(\columnwidth - 2\tabcolsep) * \real{0.2722}}
  >{\raggedright\arraybackslash}p{(\columnwidth - 2\tabcolsep) * \real{0.7278}}@{}}
\toprule\noalign{}
\begin{minipage}[b]{\linewidth}\raggedright
\textbf{Agreement}
\end{minipage} & \begin{minipage}[b]{\linewidth}\raggedright
\textbf{Key point}
\end{minipage} \\
\midrule\noalign{}
\endhead
\bottomrule\noalign{}
\endlastfoot
Ramsar (1971) & Wetlands \\
CITES (1973) & Trade in endangered species \\
Montreal (1987) & Ozone-depleting substances; most successful treaty \\
Basel (1989) & Transboundary hazardous waste \\
CBD (1992) & Biodiversity; Cartagena (biosafety, 2000), Nagoya (ABS, 2010), Kunming-Montreal GBF (2022, 30×30) \\
Kyoto (1997) & CDM, Joint Implementation, Emissions Trading; Annex-I countries\textquotesingle{} targets; 1st commitment 2008--12 \\
Paris (2015) & NDCs every 5 years; global stocktake \\
Stockholm Convention (2001) & POPs (persistent organic pollutants) \\
Minamata Convention (2013) & Mercury \\
\textbf{ISA} & International Solar Alliance --- launched by India \& France at COP21 (2015); HQ \textbf{Gurugram}; "One Sun One World One Grid" \\
\end{longtable}

\endgroup

\needspace{161pt}

\hypertarget{p1-09--7-deeper-dive--climate-science-cop-timeline-indian-programmes--disasters}{%
\section{Deeper Dive --- Climate Science, COP Timeline, Indian Programmes \& Disasters}\label{p1-09--7-deeper-dive--climate-science-cop-timeline-indian-programmes--disasters}}

\needspace{125pt}

\hypertarget{p1-09--71-greenhouse-effect--climate-change}{%
\subsection{Greenhouse Effect \& Climate Change}\label{p1-09--71-greenhouse-effect--climate-change}}

The greenhouse effect itself is natural and essential --- without it the Earth would be about 33 °C colder.
The problem is its \emph{enhancement} by human emissions, which traps extra heat.

\begin{figure}[!htbp]
\centering
\begin{adjustbox}{max width=\linewidth}
\begin{tikzpicture}[callout/.style={circle, fill=fwine, text=white, font=\scriptsize\bfseries, inner sep=1pt, minimum size=4.6mm}]
  \fill[fsage] (0,0) rectangle (12,0.5); \draw[fgink, line width=0.5pt] (0,0.5) -- (12,0.5);
  \fill[fslate!55] (0,0.5) rectangle (12,3.2);
  \draw[fgink!50, dash pattern=on 3pt off 2pt] (0,3.2) -- (12,3.2);
  \node[lbl, anchor=west] at (0.15,2.95) {atmosphere with greenhouse gases};
  \node[lbl, anchor=west] at (0.15,0.25) {Earth's surface};
  \fill[fochre] (1.2,5.0) circle (0.55);
  \foreach \a in {0,30,...,330} \draw[fochre, line width=1pt] ($(1.2,5.0)+(\a:0.7)$) -- ($(1.2,5.0)+(\a:0.9)$);
  \draw[->, fochre!80!black, line width=1.4pt] (2.0,4.5) -- (4.4,0.55);
  \node[callout] at (3.65,3.4) {1};
  \draw[->, fwine, line width=1.2pt, decorate, decoration={snake, amplitude=1.5pt, segment length=7pt, post length=2mm}] (5.6,0.55) -- (6.8,3.1);
  \node[callout] at (5.75,2.2) {2};
  \draw[->, fwine, line width=1.2pt, decorate, decoration={snake, amplitude=1.5pt, segment length=7pt, post length=2mm}] (6.8,3.15) -- (8.0,0.6);
  \node[callout] at (7.95,2.2) {3};
  \draw[->, fwine!60, line width=1.0pt, decorate, decoration={snake, amplitude=1.5pt, segment length=7pt, post length=2mm}] (9.0,0.55) -- (10.0,4.4);
  \node[callout] at (10.35,3.9) {4};
\end{tikzpicture}
\end{adjustbox}
\caption{The greenhouse effect. (1) Short-wave sunlight passes through the atmosphere and warms the
surface. (2) The warm surface radiates long-wave infrared, which greenhouse gases absorb. (3) The gases
re-emit much of it back towards the surface. (4) The rest escapes to space. More greenhouse gas means more
re-emission (3) and a warmer lower atmosphere.}
\end{figure}


\begin{itemize}
\tightlist
\item
  Pre-industrial CO₂ ≈ 280 ppm; now \textgreater{} 420 ppm. Largest anthropogenic GHG by volume: \textbf{CO₂}; methane is far more potent per molecule over 20 years.
\item
  \textbf{IPCC} (1988, UNEP + WMO) assessment reports; AR6 (2021--23): warming about 1.1 °C above pre-industrial.
\item
  Impacts: sea-level rise, glacier retreat, heatwaves, extreme rainfall, ocean acidification, coral bleaching, crop yield changes, vector-borne diseases.
\item
  \textbf{Socio-economic \& political dimensions}: climate justice; \textbf{CBDR-RC} (Common But Differentiated Responsibilities and Respective Capabilities) --- principle from Rio 1992; loss and damage fund (COP27, Sharm el-Sheikh 2022); climate finance (USD 100 billion goal; NCQG at COP29 Baku 2024).
\item
  Mitigation (reduce emissions) vs \textbf{adaptation} (adjust to impacts).
\item
  \textbf{Carbon credit} = 1 tonne CO₂-equivalent reduced. India\textquotesingle s Carbon Credit Trading Scheme (2023). \textbf{LiFE} (Lifestyle for Environment) mission launched by India at COP26.
\end{itemize}

\needspace{130pt}

\hypertarget{p1-09--72-conference-of-parties-cop--key-milestones}{%
\subsection{Conference of Parties (COP) --- Key Milestones}\label{p1-09--72-conference-of-parties-cop--key-milestones}}

\begingroup\small

\begin{longtable}[]{@{}
  >{\raggedright\arraybackslash}p{(\columnwidth - 4\tabcolsep) * \real{0.0868}}
  >{\raggedright\arraybackslash}p{(\columnwidth - 4\tabcolsep) * \real{0.2190}}
  >{\raggedright\arraybackslash}p{(\columnwidth - 4\tabcolsep) * \real{0.5270}}@{}}
\toprule\noalign{}
\begin{minipage}[b]{\linewidth}\raggedright
\textbf{COP}
\end{minipage} & \begin{minipage}[b]{\linewidth}\raggedright
\textbf{Year \& place}
\end{minipage} & \begin{minipage}[b]{\linewidth}\raggedright
\textbf{Outcome}
\end{minipage} \\
\midrule\noalign{}
\endhead
\bottomrule\noalign{}
\endlastfoot
COP1 & 1995 Berlin & Berlin Mandate \\
COP3 & 1997 Kyoto & Kyoto Protocol \\
COP8 & 2002 New Delhi & Delhi Declaration \\
COP13 & 2007 Bali & Bali Action Plan \\
COP15 & 2009 Copenhagen & Copenhagen Accord \\
COP16 & 2010 Cancun & Green Climate Fund \\
COP21 & 2015 Paris & Paris Agreement; ISA launched \\
COP26 & 2021 Glasgow & Glasgow Climate Pact; India Panchamrit, net zero 2070 \\
COP27 & 2022 Sharm el-Sheikh & Loss \& damage fund \\
COP28 & 2023 Dubai & First global stocktake; "transition away from fossil fuels" \\
COP29 & 2024 Baku & New climate finance goal (NCQG) \\
\end{longtable}

\endgroup

\needspace{130pt}

\hypertarget{p1-09--73-indian-environmental-programmes}{%
\subsection{Indian Environmental Programmes}\label{p1-09--73-indian-environmental-programmes}}

\begingroup\small

\begin{longtable}[]{@{}
  >{\raggedright\arraybackslash}p{(\columnwidth - 4\tabcolsep) * \real{0.3782}}
  >{\raggedright\arraybackslash}p{(\columnwidth - 4\tabcolsep) * \real{0.0948}}
  >{\raggedright\arraybackslash}p{(\columnwidth - 4\tabcolsep) * \real{0.5264}}@{}}
\toprule\noalign{}
\begin{minipage}[b]{\linewidth}\raggedright
\textbf{Programme}
\end{minipage} & \begin{minipage}[b]{\linewidth}\raggedright
\textbf{Year}
\end{minipage} & \begin{minipage}[b]{\linewidth}\raggedright
\textbf{Aim}
\end{minipage} \\
\midrule\noalign{}
\endhead
\bottomrule\noalign{}
\endlastfoot
Ganga Action Plan & 1985 & River cleaning → Namami Gange (2014) \\
National Clean Air Programme (NCAP) & 2019 & Cut PM by 20--30\% (revised 40\%) in non-attainment cities \\
Swachh Bharat Mission & 2014 & Sanitation, ODF \\
Pradhan Mantri Ujjwala Yojana & 2016 & LPG to poor households (reduces indoor air pollution) \\
UJALA & 2015 & LED bulbs --- energy efficiency \\
National Green Hydrogen Mission & 2023 & Green hydrogen production \\
Nagar Van Yojana & 2020 & Urban forests \\
Green India Mission & NAPCC & Increase forest/tree cover \\
Bharat Stage VI emission norms & 2020 & Leapfrogged from BS-IV \\
\end{longtable}

\endgroup

Environmental movements: \textbf{Chipko} (1973, Uttarakhand --- Sunderlal Bahuguna, Gaura Devi), \textbf{Appiko} (1983, Karnataka), \textbf{Silent Valley} (Kerala), \textbf{Narmada Bachao Andolan} (Medha Patkar), \textbf{Bishnoi} sacrifice (Khejarli, 1730, Amrita Devi).

\needspace{130pt}

\hypertarget{p1-09--74-disasters--scales--facts}{%
\subsection{Disasters --- Scales \& Facts}\label{p1-09--74-disasters--scales--facts}}

\begingroup\small

\begin{longtable}[]{@{}
  >{\raggedright\arraybackslash}p{(\columnwidth - 2\tabcolsep) * \real{0.1083}}
  >{\raggedright\arraybackslash}p{(\columnwidth - 2\tabcolsep) * \real{0.8917}}@{}}
\toprule\noalign{}
\begin{minipage}[b]{\linewidth}\raggedright
\textbf{Hazard}
\end{minipage} & \begin{minipage}[b]{\linewidth}\raggedright
\textbf{Measure / fact}
\end{minipage} \\
\midrule\noalign{}
\endhead
\bottomrule\noalign{}
\endlastfoot
Earthquake & Magnitude (Richter / moment magnitude --- logarithmic); intensity (Modified Mercalli, I--XII); India has seismic zones II--V (Zone V highest: NE, J\&K, Kutch) \\
Cyclone & IMD categories; named by WMO/ESCAP panel; Bay of Bengal more cyclone-prone than Arabian Sea \\
Tsunami & Caused by undersea earthquakes; INCOIS Hyderabad runs the warning centre \\
Flood & Most frequent disaster in India; Assam, Bihar most affected \\
Drought & Meteorological, hydrological, agricultural \\
Landslide & Himalayas and Western Ghats \\
Heat wave & Declared by IMD when max temp ≥ 40 °C (plains) and departure ≥ 4.5 °C \\
\end{longtable}

\endgroup

Disaster management cycle: \textbf{Mitigation → Preparedness → Response → Recovery}.

\needspace{95pt}

\hypertarget{p1-09--75-renewable-energy--india-snapshot}{%
\subsection{Renewable Energy --- India Snapshot}\label{p1-09--75-renewable-energy--india-snapshot}}

\begin{itemize}
\tightlist
\item
  Installed renewable capacity crossed \textbf{200 GW} (2024); solar is the largest renewable source.
\item
  Largest solar park: \textbf{Bhadla (Rajasthan)}; large hybrid park: \textbf{Khavda (Gujarat)}.
\item
  \textbf{PM-KUSUM} (solar pumps for farmers), \textbf{PM Surya Ghar Muft Bijli Yojana} (2024, rooftop solar).
\item
  Hydro: Tehri dam (tallest in India); nuclear plants: Kudankulam (largest), Tarapur (oldest).
\end{itemize}

\needspace{95pt}

\hypertarget{p1-09--previous-year-questions-pyq-pattern}{%
\section{Previous Year Questions (PYQ pattern)}\label{p1-09--previous-year-questions-pyq-pattern}}

\begin{enumerate}
\def\labelenumi{\arabic{enumi}.}
\tightlist
\item
  How many SDGs and targets? \textbf{17 goals, 169 targets}
\item
  SDG 4 is about: \textbf{Quality education}; SDG 13: \textbf{Climate action}
\item
  The Brundtland report is titled: \textbf{Our Common Future (1987)}
\item
  Minamata disease is caused by: \textbf{Mercury}
\item
  Blue baby syndrome is caused by excess: \textbf{Nitrates}
\item
  Bhopal gas tragedy gas: \textbf{Methyl isocyanate}
\item
  Which is NOT a greenhouse gas? (a) CO₂ (b) CH₄ (c) N₂ (d) N₂O --- \textbf{Ans: (c)}
\item
  Montreal Protocol relates to: \textbf{ozone-depleting substances}
\item
  Kigali Amendment aims to phase down: \textbf{HFCs}
\item
  Kyoto Protocol came into force in: \textbf{2005}
\item
  Paris Agreement aims to limit warming to: \textbf{well below 2°C, pursue 1.5°C}
\item
  ISA headquarters: \textbf{Gurugram, India}
\item
  Environment (Protection) Act enacted in: \textbf{1986}
\item
  Number of missions in NAPCC: \textbf{8}
\item
  NDMA is chaired by: \textbf{Prime Minister}
\item
  Sendai Framework period: \textbf{2015--2030}
\item
  Arrange chronologically: Stockholm (1972), Montreal (1987), Rio (1992), Kyoto (1997), Paris (2015)
\item
  Which pollutant is NOT included in India\textquotesingle s AQI? (a) PM2.5 (b) NH₃ (c) CO₂ (d) Pb --- \textbf{Ans: (c)}
\item
  India\textquotesingle s net-zero target year: \textbf{2070}
\item
  Ozone thickness is measured in: \textbf{Dobson units}
\item
  Which is a renewable resource? \textbf{Geothermal}
\item
  E-waste rules make producers responsible through: \textbf{Extended Producer Responsibility (EPR)}
\item
  CBD was opened for signature at: \textbf{Rio Earth Summit 1992}
\end{enumerate}

\textbf{More practice questions}

\begin{enumerate}
\def\labelenumi{\arabic{enumi}.}
\setcounter{enumi}{23}
\tightlist
\item
  The IPCC was established in: \textbf{1988 by UNEP and WMO}
\item
  The principle of CBDR originated at: \textbf{Rio Earth Summit 1992}
\item
  The loss and damage fund was agreed at: \textbf{COP27, Sharm el-Sheikh}
\item
  COP28 was held in: \textbf{Dubai (2023)}
\item
  National Clean Air Programme was launched in: \textbf{2019}
\item
  Chipko movement began in: \textbf{1973 (Uttarakhand)}
\item
  India\textquotesingle s highest seismic zone: \textbf{Zone V}
\item
  Tsunami early-warning centre of India is at: \textbf{INCOIS, Hyderabad}
\item
  The most frequent natural disaster in India: \textbf{floods}
\item
  The LiFE mission was announced by India at: \textbf{COP26 Glasgow}
\item
  Bhadla solar park is in: \textbf{Rajasthan}
\item
  Arrange COPs chronologically: Kyoto (COP3) → Copenhagen (COP15) → Paris (COP21) → Glasgow (COP26)
\end{enumerate}

\begin{revbox}

\begin{itemize}
\tightlist
\item
  MDG 8 (2000--15) · SDG 17/169 (2015--30)
\item
  Hg--Minamata · Cd--Itai-itai · NO₃--blue baby · As--black foot · F--fluorosis
\item
  EPA 1986 · NAPCC 2008 (8 missions) · DM Act 2005 · NGT 2010
\item
  Montreal 87 · Rio 92 · Kyoto 97 · Paris 15 · Kigali 16 · ISA 2015 (Gurugram)
\end{itemize}

\end{revbox}
